% The main theorems: the entry points of the search, the statements on the points that a search
% finds, its refusals, the delivery, the output without the exact check and the output of the
% program, the proofs in outline, the block, lag and thread lemmas, and what comparator certifies.
\chapter{The main theorems}\label{chap:theorems}

This chapter states what is proved about the search of the model, the counterpart of the library
function \verb|find_points|: which points it finds, without duplicates and in which order, which
inputs it refuses, what it delivers to a callback and without the exact check, and what the program
\verb|ratpoints-model| prints. Comparator certifies nine of these statements, the two on the
brute-force reference and the two on the fast search of Chapter~\ref{chap:overview}, thirteen in all
(Section~\ref{sec:comparator}). The proofs in outline (Section~\ref{sec:theoremsProofs}) reduce them
to what Chapters~\ref{chap:setup} to~\ref{chap:output} prove part by part: that the test of the
sieve lets every point through and that the parts compose to the search.

Throughout, $c$ is a coefficient list of degree~$d$ (Definition~\ref{def:binaryForm}), $H$ a natural
number, $s$ the options of the sieve and~$o$ the options of a search (Definition~\ref{def:Options}).
The search intervals of~$o$ are read by the values of their ends (\lname{Ratpoints.domainValues}),
except where their well-formedness is tested, which reads their bits (\lname{Ratpoints.domainBits}).
A triple \emph{lies in the search region of~$o$} when it lies in the search region of these values
and of the range of the denominators from $o.\mathrm{bLow}$ to $o.\mathrm{bHigh}$
(Definition~\ref{def:regionSpec}). The hypothesis
\[ h_S\colon\quad 0 \le o.\mathrm{sturm} \;\to\; \mathrm{SturmTheorem} \]
assumes Sturm's theorem (Definition~\ref{def:sturmTheorem}) only when the Sturm step is on, so that
a statement with~$h_S$ holds outright for a search without the Sturm step.

\section{The search}\label{sec:theoremsSearch}

The model has three entry points, one inside the next: the search with a callback and a print
channel, which the program runs; the search with a callback, the library function as a caller
without \verb|RATPOINTS_VERBOSE| sees it; and $\mathrm{findPoints}$, which collects the points in a
list. The statements speak about the last two.

\begin{definition}[The search for the points]\label{def:findPoints}
  \uses{def:Options, def:deliver, def:maxDegree, def:setupSearch, def:searchPoints, def:callback, def:libraryReport}
  \lean{Ratpoints.findPointsVerbose, Ratpoints.findPointsWith, Ratpoints.findPoints}
  \leanok
  The search with a callback and a print channel models the library function \verb|find_points|, the
  function \verb|find_points_work| that it calls and the start of \verb|find_points_work_1|
  (\csource{ratpoints.h}{436}{437}, \csource{find_points.c}{4523}{4528},
  \csource{find_points.c}{4624}{4633}, among others). It makes the refusals of the program in its
  order (Lemma~\ref{lem:setupOrder}), first those of the coefficient list, the height bound and the
  leading zeros (errors 3, 5, 4 and~8), then those of the set-up (Definition~\ref{def:setupSearch}),
  which may also return at once after the Sturm step with nothing delivered; otherwise it runs the
  search of the polynomial that the set-up hands on (Definition~\ref{def:searchPoints}). A refusal
  comes with the data of the callback as the report has left it; otherwise the search returns its
  outcome (Definition~\ref{def:Options}). The print channel prints the report of the library under
  \verb|-v| into the data of the callback (Definition~\ref{def:libraryReport}). The search with a
  callback is this search with the data dropped from a refusal and, by default, the print channel
  that prints nothing; $\mathrm{findPoints}(c, H, s, o)$ is the list of the array that it gives for
  the callback that appends each point and never stops (\lname{Ratpoints.appendPoint}), from the
  empty array, or its refusal.
\end{definition}

\begin{definition}[The search of the polynomial handed on]\label{def:searchPoints}
  \uses{def:setupUs1, def:twoAdic, def:packing, def:choiceOfModuli, def:pointsAtInfinity, def:Shape, def:correction, def:deliver, def:callback, def:libraryReport, def:reversal}
  \lean{Ratpoints.searchPoints, Ratpoints.chooseAndSearch, Ratpoints.sieving, Ratpoints.Setup.useCLong}
  \leanok
  For the polynomial that the set-up hands on, with the flags of the check (the flag of the reversal
  and the two options of the output), the range of the denominators and the intervals that the Sturm
  step leaves, with its flag \verb|sturm_empty|, the search sets up the tests on the denominators
  (\csource{find_points.c}{4900}{4914}) and the 2-adic information, and returns with nothing
  delivered when no class modulo~$64$ admits a numerator and there is no point at infinity
  (\csource{find_points.c}{4954}{4959}); it tests whether the coefficients are all \verb|long|s,
  (Definition~\ref{def:maxDegree}; \csource{find_points.c}{4814}{4823}), which chooses how the tables reduce them modulo a prime; it
  chooses the moduli
  (Definition~\ref{def:choiceOfModuli}), with nothing delivered when the curve has no points modulo
  a prime (\csource{find_points.c}{4996}{5006}); it delivers the points at infinity
  (Definition~\ref{def:pointsAtInfinity}) and returns when the callback has stopped the search or
  the Sturm step has left no interval (\csource{find_points.c}{5100}{5103},
  \csource{find_points.c}{5105}{5106}); otherwise it runs over the blocks of the denominators of the
  shape of the loop (Definition~\ref{def:Shape}) with the correction of the moduli
  (Definition~\ref{def:correction}, \csource{find_points.c}{5137}{5139}). Each point is mapped back
  to the caller's curve when the flag of the reversal is set (Definition~\ref{def:reversal}). The
  outcome has the bits of the flags, the search intervals that the Sturm step leaves (those that
  this search is given), and the numbers of the moduli of the last block that the run sifts
  (\csource{find_points.c}{5210}{5213}), $0$ after a return before the sieve.
\end{definition}

\section{The points that a search finds}\label{sec:theoremsPoints}

\begin{theorem}[Membership]\label{thm:membership}
  \uses{def:findPoints, def:IsPoint, def:regionSpec, def:sturmTheorem, def:Options, def:bruteForce}
  \lean{Ratpoints.mem_iff_isPoint_of_findPoints_eq_ok, Ratpoints.perm_filter_bruteForce_of_findPoints_eq_ok, Ratpoints.toFinset_eq_filter_points_of_findPoints_eq_ok}
  \leanok
  Assume~$h_S$ and that the exact check is on, $o.\mathrm{noCheck} = \mathrm{false}$, and let
  $\mathrm{findPoints}(c, H, s, o)$ return the list~$L$. A triple $p = (x, y, z)$ lies in~$L$ if and
  only if~$p$ is a normalized point of height at most~$H$ on $y^2 = F_c(x, z)$
  (Definition~\ref{def:IsPoint}), $y \ge 0$ if~$o.\mathrm{noY}$ is set, and~$p$ lies in the search
  region of~$o$. Indeed~$L$ is a permutation of the brute-force reference
  $\mathrm{bruteForce}(c, H)$ filtered by the last two conditions; and when neither option of the
  output is set, the set of the elements of~$L$ is $\mathrm{points}(c, H)$ filtered by the search
  region of~$o$.
\end{theorem}

\begin{proof}
  \uses{lem:findPointsList, lem:searchDens, lem:swapPoint_bijection, thm:bruteForce}
  \leanok
  By Lemma~\ref{lem:findPointsList}(3), $L$ is what the brute-force reference of the polynomial
  searched finds at infinity and over the denominators of its search up to the largest, mapped back
  after a reversal and filtered by the two conditions. Mapped back, these are a permutation of the
  points of $\mathrm{bruteForce}(c, H)$ at infinity or with a denominator up to the largest
  (Lemmas~\ref{lem:searchDens} and~\ref{lem:swapPoint_bijection}), and the search region holds no
  point with a larger one; the reference finds exactly the points (Theorem~\ref{thm:bruteForce}).
\end{proof}

\begin{theorem}[No duplicates]\label{thm:nodup}
  \uses{def:findPoints}
  \lean{Ratpoints.nodup_of_findPoints_eq_ok, Ratpoints.nodup_map_bruteForce_of_reversal_eq_ok}
  \leanok
  For all~$c$, $H$, $s$ and~$o$, if $\mathrm{findPoints}(c, H, s, o)$ returns the list~$L$, then~$L$
  has no duplicates.
\end{theorem}

\begin{proof}
  \uses{lem:findPointsList, lem:searchList, lem:searchDens, lem:swapPoint_bijection, thm:bruteForce}
  \leanok
  By Lemma~\ref{lem:findPointsList}(1), $L$ is the list of a search for a sound test. With the exact
  check on, Lemma~\ref{lem:searchList}(2) makes it a sublist of the brute-force reference of the
  polynomial searched, mapped back; the denominators of the search have no duplicates
  (Lemma~\ref{lem:searchDens}), and the reference of a reversed polynomial, mapped back, is a
  permutation of that of~$c$ (Lemma~\ref{lem:swapPoint_bijection}), which has no duplicates
  (Theorem~\ref{thm:bruteForce}). With the exact check off, Lemma~\ref{lem:searchList}(3).
\end{proof}

No duplicates needs no hypothesis at all: neither Sturm's theorem nor the exact check.

\begin{theorem}[The order of the points]\label{thm:order}
  \uses{def:findPoints, def:orderedPoints, def:regionSpec, def:sturmTheorem, def:maxDegree, def:Options}
  \lean{Ratpoints.eq_orderedPoints_of_findPoints_eq_ok, Ratpoints.pointsOfDen_eq_pointsWithDen, Ratpoints.pointsOfDen_zero_eq_pointsWithDen, Ratpoints.findPoints_sturm_neg_eq_ok}
  \leanok
  Assume~$h_S$ and that the exact check is on, $o.\mathrm{noCheck} = \mathrm{false}$. Let
  $\mathrm{findPoints}(c, H, s, o)$ return the list~$L$, let~$r$ be true when~$o$ allows the
  reversal and its search region is the whole box of~$H$, and~$B$ the largest denominator of the
  search, $o.\mathrm{bHigh}$ if $1 \le o.\mathrm{bHigh} \le H$ and~$H$ otherwise
  (Definition~\ref{def:regionSpec}). Then~$L$ is
  \[ \mathrm{orderedPoints}(\mathrm{maxDivisors}, r, c, H, B) \]
  (Definition~\ref{def:orderedPoints}) with the points left out that have $y < 0$
  when~$o.\mathrm{noY}$ is set or do not lie in the search region of~$o$. Moreover, the search with
  the Sturm step off and the other options as in~$o$ returns~$L$ too.
\end{theorem}

\begin{proof}
  \uses{lem:findPointsList, lem:swapPoint_bijection, lem:searchDens, lem:searchIntervals, def:searchRegion, lem:sturmCoverage, thm:refusals}
  \leanok
  Lemma~\ref{lem:findPointsList}(3) gives~$L$ as in the proof of Theorem~\ref{thm:membership}. The
  normalized options have the switch~$r$ and the largest denominator~$B$
  (Lemma~\ref{lem:searchIntervals}); the reversal check with~$r$ hands on
  $\mathrm{searchedPoly}(r, c)$ (Lemma~\ref{lem:swapPoint_bijection}); the denominators of its
  search from~$1$ to~$B$ are those of $\mathrm{searchDenominators}$ (Lemma~\ref{lem:searchDens});
  and over a denominator $1 \le z \le H$, or the candidate $(1, 0)$, the reference finds the points
  with the denominator~$z$, or~$0$. For the last statement, the search with the step off refuses
  only where the search with the step on refuses too (Theorem~\ref{thm:refusals}), and otherwise
  both lists are given by the first statement, which does not mention the depth.
\end{proof}

So the list is the same for every value of the options of the sieve, which change the moduli and the
work but not the points or their order; and with the exact check on, the depth of the Sturm step
does not change it either. This last statement is proved, not certified.

\section{The refusals}\label{sec:theoremsRefusals}

The refusals of the search are those of the program (Chapter~\ref{chap:setup}). Two theorems say
exactly which inputs they are, with the Sturm step on and off; so a search that refused every input
would not meet the statements of this chapter.

\begin{theorem}[The refusals]\label{thm:refusals}
  \uses{def:findPoints, def:binaryForm, def:binaryFormPoly, def:regionSpec, def:maxDegree, def:Options}
  \lean{Ratpoints.exists_findPoints_eq_error_iff_of_sturm_nonneg, Ratpoints.exists_findPoints_eq_error_iff_of_sturm_neg}
  \leanok
  Let $n + 1$ be the length of~$c$ and $N = 2 \lceil n/2 \rceil$. With the Sturm step on
  ($0 \le o.\mathrm{sturm}$), $\mathrm{findPoints}(c, H, s, o)$ refuses if and only if
  \begin{enumerate}
    \item $c$ has fewer than two or more than $\mathrm{maxDegree} + 1$ entries,
    \item $H = 0$ or $H \ge 2^{63}$,
    \item the binary form~$\Phi_c$ (Definition~\ref{def:binaryFormPoly}) is not squarefree, or
    \item the search intervals of~$o$ are not well formed (Definition~\ref{def:regionSpec}).
  \end{enumerate}
  With the Sturm step off ($o.\mathrm{sturm} < 0$), it refuses if and only if (1), (2) or (4) holds,
  or $d + 2 \le N$, or~$o$ allows the reversal, its search region is the whole box of~$H$, $d$ is
  even and $c_0 = c_1 = 0$.
\end{theorem}

\begin{proof}
  \uses{lem:setupOrder, lem:squarefree_binaryFormPoly, lem:reversal_swap, lem:chain_squarefree, lem:searchIntervals, lem:sturmStep}
  \leanok
  The search refuses in the order of Lemma~\ref{lem:setupOrder}, a natural number being a
  \verb|long| exactly when it is below~$2^{63}$. The leading zeros lower $\lceil n/2 \rceil$ exactly
  when $d + 2 \le N$, and then $Z^2 \mid \Phi_c$; the reversal check refuses exactly in the case
  stated, and then $X^2 \mid \Phi_c$ (Lemmas~\ref{lem:squarefree_binaryFormPoly}
  and~\ref{lem:reversal_swap}); malformed search intervals are refused
  (Lemma~\ref{lem:searchIntervals}); and the Sturm step, when it is on, refuses exactly when the
  remainder chain is none (Lemma~\ref{lem:sturmStep}), which after the tests before it means
  that~$\Phi_c$ is not squarefree (Lemma~\ref{lem:chain_squarefree}). As each refusal of the form
  implies that~$\Phi_c$ is not squarefree, the disjunctions come out as stated.
\end{proof}

Without the Sturm step the search refuses only some of the forms that are not squarefree, as
\ratpoints{} does under \verb|-s|, and searches the others; the statements on the points hold for these
searches too.

\section{The delivery and the options of the output}\label{sec:theoremsDelivery}

The library delivers each point to the callback where it finds it and stops when a call sets the
flag \verb|*quit| (Chapter~\ref{chap:output}). By the protocol this comes to a fold of the list of
$\mathrm{findPoints}$ (Definition~\ref{def:deliver}), so that the statements on the list carry over
to every callback.

\begin{theorem}[The protocol of the delivery]\label{thm:protocol}
  \uses{def:findPoints, def:deliver, def:Options}
  \lean{Ratpoints.map_findPointsWith_eq_map_findPoints, Ratpoints.info_total_eq_of_findPointsWith_eq_ok, Ratpoints.exists_findPointsWith_eq_ok_of_findPoints_eq_ok}
  \leanok
  For all~$c$, $H$, $s$ and~$o$, every callback and its data: the search with the callback, with the
  print channel that prints nothing, refuses exactly when $\mathrm{findPoints}(c, H, s, o)$ refuses,
  with the same refusal; and when $\mathrm{findPoints}(c, H, s, o)$ returns the list~$L$, the search
  with the callback returns an outcome whose data and count are those of the delivery of~$L$ to the
  callback from its data, the count~$0$ and the flag of the stop clear. So the calls of the callback
  are those for the points of~$L$, in their order, up to and including the first call that stops the
  search, and the count is the sum of the weights that these calls return.
\end{theorem}

\begin{proof}
  \uses{lem:setupOrder, lem:searchList, lem:searchIntervals, lem:sturmStep}
  \leanok
  Both searches make the same refusals, which do not depend on the callback, and otherwise search
  the same polynomial over the same range of the denominators and the intervals that the Sturm step
  leaves, or both return at once after the Sturm step (Lemma~\ref{lem:setupOrder}). The intervals
  are well formed and the range lies in $1, \dots, H$ (Lemmas~\ref{lem:searchIntervals}
  and~\ref{lem:sturmStep}), so that Lemma~\ref{lem:searchList}(1) gives the data and the count of
  the search with the callback as those of the delivery of the list of the search with the appending
  callback. The proof rests on the contract of the search, with one test of the sieve for every
  callback, not on an argument about the control flow.
\end{proof}

The protocol speaks of the data and the count only. The rest of the outcome is described in
Section~\ref{sec:outputBack}; the numbers of the moduli used depend on where the callback stops the
search. The report under \verb|-v|, which comes through the print channel, is compared with that of
the C program and not proved (Chapter~\ref{chap:ccode}).

Under \verb|RATPOINTS_NO_CHECK| the library delivers, each with $y = 0$ and mapped back after a
reversal, the candidate at infinity of the polynomial searched when that polynomial has points at
infinity (odd degree or a square leading coefficient), and every numerator that passes the sieve,
the test for common factors and the third stage, whether or not it gives a point. Which candidates
pass depends on the moduli and is not specified; the theorem below says that each is a triple
$(x, 0, z)$ with $(x, z)$ normalized, of height at most~$H$ and in the search region, and that they
include the points, in the order in which the search with the exact check finds them.

\begin{theorem}[The output without the exact check]\label{thm:noCheck}
  \uses{def:findPoints, def:IsPoint, def:regionSpec, def:sturmTheorem, def:Options}
  \lean{Ratpoints.isNormalizedPair_of_mem_of_findPoints_eq_ok_of_noCheck, Ratpoints.mem_of_isPoint_of_findPoints_eq_ok_of_noCheck, Ratpoints.map_sublist_of_findPoints_eq_ok_of_noCheck}
  \leanok
  Let the exact check be off, $o.\mathrm{noCheck} = \mathrm{true}$, and let the search
  $\mathrm{findPoints}(c, H, s, o)$ return the list~$L$.
  \begin{enumerate}
    \item Every element $(x, y, z)$ of~$L$ has $y = 0$ and the pair $(x, z)$ normalized and of
      height at most~$H$ (Definition~\ref{def:IsPoint}), and lies in the search region of~$o$.
    \item Assume~$h_S$. For every normalized point $(x, y, z)$ of height at most~$H$ in the search
      region of~$o$, the triple $(x, 0, z)$ lies in~$L$.
    \item If the search with the same~$s$ and~$o$ with $\mathrm{noCheck}$ cleared and~$\mathrm{noY}$
      set returns the list~$L'$, then~$L'$ with each~$y$ replaced by~$0$ is a sublist of~$L$.
  \end{enumerate}
\end{theorem}

\begin{proof}
  \uses{lem:findPointsList, lem:searchList, def:searchRegion, lem:searchIntervals, lem:reversal_swap, thm:membership}
  \leanok
  By Lemma~\ref{lem:findPointsList}(1), $L$ and~$L'$ are the lists of searches with the same sound
  test over the same region, which differ only in the flags of the check. Item~(1): the pair by
  Lemma~\ref{lem:searchList}(3); the region, without a reversal, by Lemma~\ref{lem:searchList}(4),
  since the region that the search uses lies in that of~$o$ (Definition~\ref{def:searchRegion},
  Lemma~\ref{lem:searchIntervals}), and after a reversal, which comes only for the whole box
  (Lemma~\ref{lem:reversal_swap}), since every normalized pair of height at most~$H$ lies in it.
  Item~(3) is Lemma~\ref{lem:searchList}(5). Item~(2): $(x, |y|, z)$ lies in~$L'$ by
  Theorem~\ref{thm:membership}, and item~(3) puts $(x, 0, z)$ into~$L$.
\end{proof}

Items~(1) and~(3) do not need Sturm's theorem. The two searches of item~(3) choose the same moduli
and make the same corrections, which depend neither on the callback nor on the options of the output
(Chapter~\ref{chap:moduli}).

\section{What the program prints}\label{sec:theoremsProgram}

The program \verb|ratpoints-model| runs the search with the callback of \verb|main.c|, which prints
the points (Chapter~\ref{chap:output}). For a command line that it reads as~$c$, $H$, $s$, $o$ and
the settings of the output, it searches with the options~$s'$ of~$s$ with the array size that
\verb|main.c| sets and the options~$o'$ of~$o$ with the search intervals of the command line joined
(Definition~\ref{def:cliNumbers}).

\begin{theorem}[What the program prints]\label{thm:cliRun}
  \uses{def:cliRun, def:cliProcess, def:cliInput, def:cliNumbers, def:findPoints, def:deliver, def:IsPoint, def:regionSpec, def:sturmTheorem}
  \lean{Ratpoints.Cli.run_eq_of_findPointsWith_eq_ok, Ratpoints.Cli.total_eq_of_findPointsWith_eq_ok, Ratpoints.Cli.run_eq_of_findPoints_eq_ok, Ratpoints.Cli.run_eq_of_onePoint, Ratpoints.Cli.inRegion_iff_of_input_eq_ok, Ratpoints.Cli.intervalsWF_joinIntervals_of_input_eq_ok, Ratpoints.Cli.mem_filter_keeps_iff_of_input_eq_ok}
  \leanok
  Let the program read a command line as above, and call the points of a list that the callback of
  the program keeps (Definition~\ref{def:cliProcess}), under \verb|-1| only the first of them, its
  \emph{counted points} (\lname{Ratpoints.Cli.countedPoints}).
  \begin{enumerate}
    \item Without \verb|-v|: if $\mathrm{findPoints}(c, H, s', o')$ returns the list~$L$ and the
      search with the callback of the program the outcome~$r$, the program prints the head of its
      report, the string before the points, the counted points of~$L$ (none under \verb|-z|) in the
      format, separated by the string between the points, the string after the points and the end of
      the report for~$r$, with the status~$0$; the count of~$r$ is the number of the counted points
      of~$L$.
    \item Under \verb|-q|, which clears \verb|-v|, the program prints, whenever
      $\mathrm{findPoints}(c, H, s', o')$ returns~$L$, the same without the head and the end of the
      report; under \verb|-q -1| without \verb|-z| the one point printed is the first point of~$L$
      that the callback keeps.
    \item The callback keeps a point unless it lies at infinity under \verb|-i|, or $z \ne 0$, the
      command line gives search intervals and~$x/z$ does not lie in their exact intervals. When
      every argument of the command line has fewer than~$2^{56}$ characters, the joined intervals
      are well formed (Definition~\ref{def:regionSpec}), so that the library does not refuse them.
    \item Assume~$h_S$ and that the exact check is on. If $\mathrm{findPoints}(c, H, s', o')$
      returns~$L$, a triple lies in~$L$ and is kept by the callback exactly when it is a normalized
      point $(x, y, z)$ of height at most~$H$ (Definition~\ref{def:IsPoint}) with $y \ge 0$ under
      \verb|-y| and $z \ne 0$ under \verb|-i|, and, for $z \ne 0$, with~$z$ in the range of the
      denominators of \verb|-dl| and \verb|-du| and~$x/z$ in the exact intervals of the command line
      when there are such intervals.
  \end{enumerate}
\end{theorem}

\begin{proof}
  \uses{thm:protocol, thm:membership, lem:cliNumbers}
  \leanok
  The delivery of any list to the callback of the program prints its counted points in turn and
  counts them, stopping under \verb|-1| at the first point kept
  (\lname{Ratpoints.Cli.deliver\_process}); by the protocol (Theorem~\ref{thm:protocol}) the search
  with the callback gives the data and the count of the delivery of~$L$; without \verb|-v|, which
  \verb|-q| clears (\lname{Ratpoints.Cli.verbose\_eq\_false\_of\_input\_eq\_ok}), the print channel
  prints nothing. Items~(3) and~(4) rest on an invariant of the reading of the intervals
  (\lname{Ratpoints.Cli.Intervals.ReadsWF}): the exact ends are those of increasing and disjoint
  intervals, and the doubles of each interval hold it for $|x| \le H$. For arguments of fewer
  than~$2^{56}$ characters (\verb|scan_end| stops accumulating the digits of an exponent at
  \verb|LONG_MAX/20|) the doubles of the lower ends after the first are the largest doubles at most
  their exact ends (Lemma~\ref{lem:cliNumbers}), so that the join makes intervals that are
  increasing and strictly apart. Item~(4) is then membership (Theorem~\ref{thm:membership}) on the
  joined intervals with the test of the callback.
\end{proof}

\section{The search as a filter of the reference}\label{sec:theoremsProofs}

The theorems come from one description of what a search delivers, in terms of a test of the sieve of
which the proofs use two properties: it lets the numerator of every point through, and it depends
neither on the callback nor on the options of the output. In the search it is the 2-adic word of the
class of the denominator (Chapter~\ref{chap:denominators}), the rows of the moduli of its block
(Chapters~\ref{chap:moduli} and~\ref{chap:rows}) and the test for common factors
(Chapter~\ref{chap:sieve}); the exact check then decides, so that the list of a search is a filter
of the brute-force reference, up to the order and the region.

\begin{definition}[The list of a search]\label{def:searchList}
  \uses{def:pointsAtInfinity, def:checkNumWith, def:Shape, def:regionSpec, def:bruteForce}
  \lean{Ratpoints.searchList, Ratpoints.IsSoundTest}
  \leanok
  For the coefficient list~$c$ of its degree, $H$, the denominators from $b_{\mathrm{low}}$ to
  $b_{\mathrm{high}}$, intervals~$D$ of values, a test~$S(b, a)$ of the numerators~$a$ over the
  denominators~$b$, and the flags of the check: the points at infinity with the flags
  (Definition~\ref{def:pointsAtInfinity}), followed, for the denominators~$b$ of the search from
  $b_{\mathrm{low}}$ to $b_{\mathrm{high}}$ in the order of its loop (Definition~\ref{def:Shape}),
  by the points with the flags (Definition~\ref{def:checkNumWith}) of the numerators
  $a = -H, \dots, H$ with~$a/b$ in one of the intervals of~$D$ and~$S(b, a)$. The test is
  \emph{sound} when, over each denominator~$b$ of the search, it holds for every numerator
  $a \in [-H, H]$ over which the brute-force reference finds a point $(a, y, b)$.
\end{definition}

\begin{lemma}[The list of a sound test]\label{lem:searchList}
  \uses{def:searchList, def:searchPoints, def:regionSpec, def:deliver, def:bruteForce, def:IsPoint, def:checkNumWith, lem:choiceOfModuli, lem:runRows}
  \lean{Ratpoints.exists_isSoundTest_run_searchPoints, Ratpoints.run_searchPoints_eq_deliver, Ratpoints.exists_isSoundTest_run_sieveRun, Ratpoints.Sieve.IsSoundRun, Ratpoints.searchList_of_noCheck_eq_false, Ratpoints.exists_eq_of_mem_searchList_of_noCheck, Ratpoints.nodup_searchList_of_noCheck, Ratpoints.mem_region_of_mem_searchList, Ratpoints.map_searchList_sublist}
  \leanok
  For the list of a search of~$c$ with the denominators from $b_{\mathrm{low}}$ to
  $b_{\mathrm{high}}$, intervals~$D$, a test~$S$ and the flags of the check:
  \begin{enumerate}
    \item Let $1 \le H < 2^{63}$, $1 \le b_{\mathrm{low}}$, $b_{\mathrm{high}} \le H$, and let the
      search intervals be well formed. For the normalized options of the sieve and either value of
      \verb|sturm_empty| there is a sound test~$S$ such that the search of~$c$ of its degree
      (Definition~\ref{def:searchPoints}), with the print channel that prints nothing, gives, for
      every callback, every setting of the flags of the check and every delivery whose flag of the
      stop is clear, the data and the count of the delivery of the list of the search for~$S$, the
      values of the search intervals (none when \verb|sturm_empty| is set) and these flags. So the
      search with any callback gives the data and the count of the delivery of the list that the
      search with the appending callback gives. The run over the blocks has the corresponding
      contract, for sound arguments (\lname{Ratpoints.Sieve.IsSoundRun}), a well-formed first
      decision of the choice and the bounds of well-formed search intervals.
    \item If~$S$ is sound, $d \ge 1$ and the exact check is on, the list is what the brute-force
      reference finds over the candidate $(1, 0)$, followed by what it finds over the denominators
      of the search and the numerators $a \in [-H, H]$ with~$a/b$ in~$D$, each point mapped back
      when the flag of the reversal is set, without those with $y < 0$ under~$\mathrm{noY}$.
    \item With the exact check off, $1 \le H$, $1 \le b_{\mathrm{low}}$ and
      $b_{\mathrm{high}} \le H$, whatever~$D$ and~$S$, every element is $(x, 0, z)$ with $(x, z)$
      normalized and of height at most~$H$, and the list has no duplicates.
    \item Without the flag of the reversal and for $1 \le b_{\mathrm{low}}$, every element has
      $z = 0$, or $b_{\mathrm{low}} \le z \le b_{\mathrm{high}}$ and~$x/z$ in~$D$.
    \item For the same~$D$, $S$ and flag of the reversal, the list with the exact check on
      and~$\mathrm{noY}$ set, each~$y$ replaced by~$0$, is a sublist of the list with the exact
      check off, for either value of~$\mathrm{noY}$.
  \end{enumerate}
\end{lemma}

\begin{proof}
  \uses{lem:twoAdic_point, lem:choiceOfModuli, lem:exitByPrime, def:pointsAtInfinity, thm:blocks, lem:numeratorRanges, lem:shapes, lem:rowsAdmitPoints, lem:runRows, def:sift, lem:checkNumWith, lem:deliverList, lem:swapPoint_bijection}
  \leanok
  (1): by cases on the 2-adic information and the result of the choice of the moduli, which depend
  neither on the callback nor on the flags; the test is chosen after them. At either exit the curve
  has no points (Lemmas~\ref{lem:twoAdic_point}, \ref{lem:choiceOfModuli}
  and~\ref{lem:exitByPrime}), and the test that lets nothing through is sound. Otherwise the result
  of the choice is what the run needs (Lemma~\ref{lem:choiceOfModuli}), and after the points at
  infinity the run delivers, by the block lemma (Theorem~\ref{thm:blocks}) and the contract of one
  denominator (Definition~\ref{def:sift}), the points with the flags of the numerators in the ranges
  of the intervals (Lemma~\ref{lem:numeratorRanges}) that pass the class word of~$b$ and survive the
  entries of its block. This test is sound, since the shape lets the denominator of every point
  through in the order of the search (Lemmas~\ref{lem:shapes} and~\ref{lem:searchDens}), the class
  word admits the numerator of every point (Lemma~\ref{lem:twoAdic_point}), and every entry is sound
  (Lemma~\ref{lem:rowsAdmitPoints}). (2): a sound test leaves out only numerators over which the
  reference finds nothing; (2) and~(5) then hold numerator by numerator and at infinity
  (Lemma~\ref{lem:checkNumWith}). (3): the survivors are $(a, 0, b)$ with $\gcd(a, b) = 1$,
  $|a| \le H$ and~$b$ in the range, distinct over distinct denominators
  (Lemma~\ref{lem:searchDens}), and the map of the reversal is injective on normalized pairs
  (Lemma~\ref{lem:swapPoint_bijection}). (4) is read off the definition.
\end{proof}

\begin{lemma}[The list of a search for the points]\label{lem:findPointsList}
  \uses{def:findPoints, def:searchList, def:searchRegion, lem:setupOrder, def:regionSpec, def:sturmTheorem, def:bruteForce, lem:searchDens, def:remainderChain}
  \lean{Ratpoints.exists_searchList_of_findPoints_eq_ok, Ratpoints.exists_eq_of_findPoints_eq_ok, Ratpoints.exists_eq_filter_inIntervals_of_findPoints_eq_ok, Ratpoints.exists_eq_filter_inRegion_of_findPoints_eq_ok, Ratpoints.elim_searchRegion_pointsOfDen_filter, Ratpoints.map_append_flatMap_filter_eq_filter_inRegion}
  \leanok
  Let $\mathrm{findPoints}(c, H, s, o)$ return the list~$L$, let~$c'$ and~$\mathrm{rev}$ be the list
  and the flag that the reversal check of the normalized options hands on, and let the normalized
  range of the denominators run from $b'_{\mathrm{low}}$ to $b'_{\mathrm{high}}$.
  \begin{enumerate}
    \item The predicate \lname{Ratpoints.IsSearchedPoly} of what a successful search guarantees
      after its refusals (Lemma~\ref{lem:setupOrder}) holds for~$c'$ and~$\mathrm{rev}$, and there
      is a sound test~$S$ over that range such that~$L$ is the list of the search of~$c'$ for~$S$
      and the region that the search uses (Definition~\ref{def:searchRegion}), with the
      flag~$\mathrm{rev}$ and the options of the output of~$o$, or the empty list when that region
      is none; for every choice of the two options of the output, with the other options unchanged,
      $\mathrm{findPoints}$ returns the list for the same~$S$ and region with that choice.
    \item With the exact check on, $L$ is what the brute-force reference of~$c'$ finds over the
      candidate $(1, 0)$ and over the denominators of its search from $b'_{\mathrm{low}}$ to
      $b'_{\mathrm{high}}$ with the numerators in the region that the search uses, mapped back after
      a reversal, without the points with $y < 0$ under~$\mathrm{noY}$.
    \item If moreover~$h_S$ holds, $L$ is what the reference of~$c'$ finds over the candidate
      $(1, 0)$ and over the denominators of its search from~$1$ to $b'_{\mathrm{high}}$ with all the
      numerators $-H, \dots, H$, mapped back after a reversal, without the points with $y < 0$
      under~$\mathrm{noY}$ and those outside the search region of~$o$.
  \end{enumerate}
\end{lemma}

\begin{proof}
  \uses{lem:searchList, lem:setupOrder, lem:searchIntervals, lem:sturmStep, def:searchRegion, lem:searchDens, lem:reversal_swap}
  \leanok
  (1): after the refusals the search of~$c'$ runs over the normalized range, which lies in
  $1, \dots, H$, and over well-formed intervals (Lemmas~\ref{lem:setupOrder},
  \ref{lem:searchIntervals} and~\ref{lem:sturmStep}), so that Lemma~\ref{lem:searchList}(1) applies
  with the appending callback; the options of the output enter only the flags of the check, which
  the test precedes. (2) follows by Lemma~\ref{lem:searchList}(2). (3): under~$h_S$, the region that
  the search uses holds, among the numerators of the points of a denominator, exactly those
  with~$a/b$ in the caller's intervals, and when the step returns at once the polynomial searched
  has no point at infinity (Definition~\ref{def:searchRegion}); the denominators from
  $b'_{\mathrm{low}}$ on are those from~$1$ that are at least $b'_{\mathrm{low}}$
  (Lemma~\ref{lem:searchDens}); and a reversal comes only for the whole box
  (Lemma~\ref{lem:reversal_swap}).
\end{proof}

\section{The blocks, the lag and the threads}\label{sec:theoremsBlocks}

The run sifts the denominators in blocks, each with the moduli of the decision of the correction in
force for it, and the correction may append a decision after each block (Chapter~\ref{chap:moduli}).
A decision is the first block for which it holds, with the ends~$s_2$ and~$s_3$ of the second and
third stages from that block on; the decision in force for block~$i$ is the last whose first block
is at most~$i$ (Definition~\ref{def:blocks}). The block lemma is a step of the proofs above
(Lemma~\ref{lem:searchList}). The lag lemma and the thread lemma are not used by them: they state
the property on which the pool of threads of the C code relies, so that the output of the model,
which sifts the blocks in turn, is also that of \ratpoints{} with threads. None of the three is a
statement of the challenge.

\begin{theorem}[The block lemma]\label{thm:blocks}
  \uses{def:blocks, lem:blocks, def:correction, lem:correction, lem:choiceOfModuli, lem:runRows, def:deliver, def:checkNumWith, def:exactBounds}
  \lean{Ratpoints.Sieve.blockPoints, Ratpoints.Sieve.SieveRunLoopInvariant, Ratpoints.Sieve.sieveRun_spec, Ratpoints.Sieve.exists_run_sieveRun_eq_deliver}
  \leanok
  Let the run satisfy the hypotheses of Lemma~\ref{lem:blocks} other than those on block~$i$ and its
  decision, the sieve entries having what \lname{Ratpoints.Sieve.IsRunEntries} asks (their tables of
  inverses, the reciprocals of their moduli in their order, the row offsets of the classes with a
  pattern for all of them, and prime entries where they count no factor; they need not be sound),
  let the word of the classes mark those with a nonzero pattern, and let the decision of the choice
  be well formed. Then the run delivers the points of its blocks in order, block~$i$ those of its
  denominators that pass the tests of the shape, each with the first~$s_3$ entries of the decision
  of block~$i$, which depends neither on the callback nor on the flags of the output. So, when the
  loop of the shape visits no denominator twice, there is for each denominator~$b$ a
  number~$s_3(b)$, independent of the callback and the flags, such that the run delivers, over the
  denominators of the loop that pass its tests and in its order, the points with the flags of the
  numerators in the ranges of~$b$ that pass the class word of~$b$ and survive the first~$s_3(b)$
  entries.
\end{theorem}

\begin{proof}
  \uses{lem:blocks, lem:correction, lem:deliverList}
  \leanok
  The invariant of the loop over the blocks (\lname{Ratpoints.Sieve.SieveRunLoopInvariant}): the
  delivery is that of the points of the blocks done, and until a stop the state of the correction is
  its closed form after their folds and the store fits all the entries. Each block delivers its
  points by Lemma~\ref{lem:blocks}, with a well-formed decision by Lemma~\ref{lem:correction}; a
  stop in a block ends the run before any decision that its fold appends is read, and the delivery
  after a stop is that of the whole list (Lemma~\ref{lem:deliverList}). The form per denominator
  follows, as the blocks, concatenated, visit the denominators of the loop in its order
  (Lemma~\ref{lem:blocks}).
\end{proof}

\begin{theorem}[The lag lemma]\label{thm:lag}
  \uses{def:blocks, def:correction, lem:correction}
  \lean{Ratpoints.Sieve.Decision.InForce.unique, Ratpoints.Sieve.getElem!_findDecision_correctionAfter}
  \leanok
  For the arguments of a run and every block~$i$, the decision in force for block~$i$ among the
  decisions after the folds of the first~$k$ blocks is the decision of block~$i$ for every
  $k \ge i - \mathrm{adaptLag}$: block~$i$ reads only decisions made at the folds of blocks at least
  $\mathrm{adaptLag} + 1$ before it.
\end{theorem}

\begin{proof}
  \uses{lem:correction, def:exactBounds}
  \leanok
  The decision in force for a block, the last whose first block is at most it, is unique, and
  decisions appended for later blocks change neither it nor its index. After~$k$ folds the decisions
  are those after any $m \le k$ folds followed by decisions for the blocks from
  $m + \mathrm{adaptLag} + 1$ on (Lemma~\ref{lem:correction}); take $m = i - \mathrm{adaptLag}$.
\end{proof}

\begin{theorem}[The thread lemma]\label{thm:threads}
  \uses{thm:blocks, thm:lag, lem:blocks, lem:correction, lem:runRows, def:deliver, def:exactBounds}
  \lean{Ratpoints.Sieve.deliverSlots, Ratpoints.Sieve.deliverSlots_spec, Ratpoints.Sieve.toList_info_run_runBlock_appendPoint, Ratpoints.Sieve.run_deliverSlots_eq_fst_run_sieveRun}
  \leanok
  Under the hypotheses of Theorem~\ref{thm:blocks}, let each block~$i$ be sifted alone with the
  appending callback into a slot, from the empty array, with the decisions after the folds of any
  $k_i \ge i - \mathrm{adaptLag}$ blocks and a store of its own that fits all the entries. Then
  delivering the slots in the order of the blocks to a callback whose flag is clear, with a return
  after the slot in which a call stops the search, gives what the run delivers.
\end{theorem}

\begin{proof}
  \uses{thm:lag, lem:blocks, lem:deliverList, thm:blocks, def:checkNumWith}
  \leanok
  By the lag lemma (Theorem~\ref{thm:lag}) the decision of each slot is that of its block, and the
  points of a block do not depend on the store that it starts from as long as the store fits
  (Lemma~\ref{lem:blocks}); so slot~$i$ holds the points of block~$i$, and the delivery of the slots
  in order is the delivery of their concatenation (Lemma~\ref{lem:deliverList}), which is what the
  run delivers (Theorem~\ref{thm:blocks}).
\end{proof}

The pool of threads of the C code meets these hypotheses: a thread starts block~$i$ only once the
blocks before block $i - L$ are folded, $L$ being \verb|RP_ADAPT_LAG|
(\csource{find_points.c}{4292}{4296}); the tables of rows that the threads share only gain complete
rows, published once built (\csource{init.c}{180}{180}, \csource{init.c}{241}{241},
\csource{init.c}{391}{391}, \csource{init.c}{443}{443}), so that they fit all the entries whenever a
thread reads them (\csource{find_points.c}{3160}{3160}, \csource{find_points.c}{3533}{3535}); and
the calling thread delivers the slots in the order of the blocks. The threads themselves stay
outside the model (Chapter~\ref{chap:ccode}), and \verb|ratpoints-model| accepts \verb|-t| and
ignores it.

\section{What comparator certifies}\label{sec:comparator}

Comparator, a checker for Lean proofs built into Lake, takes a challenge, a file that states
theorems with \verb|sorry|, and a solution, a file that proves them, and checks the solution
independently of the build of the repository. The files are in the folder \verb|comparator|: the
challenge \verb|Challenge.lean|, the solution \verb|Solution.lean|, and \verb|challenges.json|,
which lists the thirteen statements and the permitted axioms; \verb|scripts/run-comparator.sh| runs
the check.

The challenge may import only Mathlib, the model and the fast search, which the layout check
enforces. It restates verbatim the definitions of the four files of the specification with their
decidability instances (Chapter~\ref{chap:spec}), all but the finite set $\mathrm{points}(c, H)$
(Section~\ref{sec:specBruteForce}), and imports exactly the Mathlib modules among the public imports
of these files, so that the restated definitions elaborate to the same terms as there. It imports
the whole model (\verb|Ratpoints.Model|), of which the statements mention the search with and
without a callback and the brute-force reference; the options, the options of the sieve and the
doubles of the search intervals; the outcome, the answer of a callback, the delivery and its
reference; and the constants $\mathrm{maxDegree}$ and $\mathrm{maxDivisors}$. It imports the fast
search (\verb|Ratpoints.Fast|) too, of which the statements mention
\verb|Ratpoints.FastMask.fastPoints| and its parameters \verb|Ratpoints.Fast.Tuning|. A definition
of the model or of the fast search elaborates without Mathlib, to other instance terms, and is never
restated: they count as the code that the executables run, while the specification is reproduced, so
that a library that specified something else would be rejected.

Eleven of the statements are statements of this blueprint, each named as its theorem in the library
with \verb|challenge_| in place of \verb|Ratpoints.| (the challenge spells out the maps of the
search intervals to their values and bits, which the library abbreviates): the two statements of the
brute-force reference (Theorem~\ref{thm:bruteForce}), the first statements of membership and of the
order (Theorems~\ref{thm:membership} and~\ref{thm:order}), no duplicates (Theorem~\ref{thm:nodup}),
the two refusal theorems (Theorem~\ref{thm:refusals}), the three statements without the exact check
(Theorem~\ref{thm:noCheck}) and the first statement of the protocol (Theorem~\ref{thm:protocol}).
They hold for every coefficient list, height bound, options of a search and options of the sieve, so
for every choice of the moduli and every correction that the search makes; Sturm's theorem enters as
the hypothesis~$h_S$ of membership, the order and item~(2) without the exact check. The other two
are the theorems of the fast search, \verb|Ratpoints.FastMask.mem_fastPoints_iff|, which assumes
Sturm's theorem too, and \verb|Ratpoints.FastMask.nodup_fastPoints| (Chapter~\ref{chap:overview}). The solution proves each
statement by its theorem in the library.

\verb|lake comparator| builds and exports both files in a sandbox, without loading the \verb|.olean|
files of the project into its own process, and accepts the solution when each of its theorems states
the same as the theorem of the challenge, comparing the full bodies, not only the types, of every
definition that the statement reaches; when the proofs use only the permitted axioms,
\verb|propext|, \verb|Quot.sound| and \verb|Classical.choice|, so that a proof by
\verb|native_decide| or \verb|bv_decide| would fail (the layout check forbids both in the project);
and when the kernel accepts them. With \verb|--paranoid| it checks the solution with six kernels:
Lean's kernel in its default and in its paranoid mode, \verb|lean4lean|, \verb|nanoda|,
\verb|con-leche| and \verb|con-ron| (the flag \verb|enable_nanoda| of \verb|challenges.json| does
not switch \verb|nanoda| off under \verb|--paranoid|). On 2026-10-31,
\verb|lake comparator --paranoid| accepted the thirteen statements with all six.

What one has to read, then, is the challenge: the specification as restated there and, of the model,
the types that the statements mention and the reference of the delivery, a fold of a few lines.
Comparator does not certify that the executable \verb|ratpoints-model|, which imports only the model
and is compiled by the compiler of Lean, runs the model: the layout check keeps \verb|partial|,
\verb|unsafe|, \verb|implemented_by|, \verb|extern| and experimental options out of the model's own
code (Section~\ref{sec:ccodeCI}), while the primitives of Lean core that the compiler implements
natively, on machine words, arrays, thunks and \verb|Float|, remain trusted
(Section~\ref{sec:overviewTrust}); no proof reads the value of a \verb|Float|
(Remark~\ref{rem:moduliFloats}). Nor does comparator say anything about the C code
(Chapter~\ref{chap:ccode}), or about the theorems outside the challenge (the forms of membership
with the permutation and the set of the points, the depth of the Sturm step, the program theorem,
the block, lag and thread lemmas): they are proved in the same library, but one has to read their
Lean statements to know what they say.
